\documentclass[11pt]{article}
\usepackage[letterpaper,margin=1in]{geometry}
\usepackage{mathtools}
\usepackage{hyperref}

\usepackage[backend=biber,
  maxbibnames=99,
  natbib=true,
  style=ext-numeric-comp,
%  autocite=numeric, %superscript,
  sortcites=true,
  sorting=nyt]{biblatex}
\DeclareFieldFormat % Titles in sentence case
  [article,inbook,incollection,inproceedings,patent,thesis,unpublished]
  {titlecase:title}{\MakeSentenceCase*{#1}}
\renewbibmacro{in:}{ % Don't use "In:" before journal names.  
  \ifentrytype{article}{}{\printtext{\bibstring{in}\intitlepunct}}}
% date after authors
\renewbibmacro*{date+extradate}{}
\renewbibmacro*{issue+date}{}
\renewbibmacro*{date+extradate}{%
  \printtext[parens]{%
    \printdate
  }%
}
\renewbibmacro*{author}{%
  \ifuseauthor
    {\printnames{author}%
     \setunit{\addspace}%
     \usebibmacro{date+extradate}}
    {}%
}

\addbibresource{defs.bib}
\addbibresource{arr.bib}
\addbibresource{molrec.bib}

\renewcommand{\topfraction}{.9}
\renewcommand{\bottomfraction}{.9}
\renewcommand{\textfraction}{.1}
\renewcommand{\floatpagefraction}{.8}
\renewcommand{\dbltopfraction}{.9}
\renewcommand{\dblfloatpagefraction}{.9}
\setcounter{topnumber}{2}
\setcounter{bottomnumber}{2}
\setcounter{totalnumber}{5}
\setcounter{dbltopnumber}{2}
\setcounter{secnumdepth}{0}
\begin{document}
\title{The Misapportionment of Human Diversity}
\author{Alan R. Rogers}
\maketitle

Consider a set of $K$ loci, which segregate in two populations of
equal size. At each locus there are two alleles, called ``plus'' and
``minus'' for reasons explained below. The frequency of each plus
allele is $p_1 = 1/2 + \delta$ in population~1, $p_2 = 1/2 - \delta$
in population~2, and $1/2$ in the population as a whole.  The
between-group component of variance in allele frequencies is $B = (p_1
- 1/2)^2/2 + (p_2 - 1/2)^2/2 = \delta^2$. The within-group component
is $W = 2p_1(1-p_1)/2 + 2p_2(1-p_2)/2 = 1/2 - 2\delta^2$. This implies
that
\begin{eqnarray*}
  F_{ST} &=& B/(B+W)\\
  &=& \frac{\delta^2}{\delta^2 + 1/2 - 2\delta^2}\\
  &\approx& 2\delta^2
\end{eqnarray*}
where the approximation assumes that $\delta^2$ is small. $F_{ST}$
measures the between-group variance in allele frequencies as a
fraction of total genetic variance. Although Lewontin did not use
$F_{ST}$, we can paraphrase his main result by saying that $F_{ST}$
among the major human populations is small---that differences between
these populations account for only a small fraction of human genetic
variance. To make a point, I will assume that $F_{ST}$ is even smaller
than the values typically estimated from human data.

Given that $F_{ST}$ is small, what does this implies about group
differences in a polygenic character controlled by these same loci? I
assume that, at each locus, the plus allele adds one unit to character
value, and the minus allele subtracts one unit. The environment is the
same in both groups and therefore makes no contribution to
between-group variance.  With this setup, each locus contributes
$4\delta$ to the mean character value in population~1 and $-4\delta$
to that in population~2. Each locus thus contributes $8\delta$ to $D$,
the difference between populations and $16\delta^2$ to $V_B$, the
between-population component of variance. Summing across loci,
\begin{eqnarray*}
D &=& 8K\delta\\
V_B &=& 16 K \delta^2
\end{eqnarray*}
The variance here is a sum of per-locus contributions, because I have
assumed linkage equilibrium between all pairs of loci.

Environmental effects become important as we turn to the within-group
component of variance. The heritability within each population is by
definition $h^2 = V_G/(V_G + V_E)$, where $V_G$ and $V_E$ are the
additive genetic and environmental components of variance. The total
variance within groups can thus be written as $V_G + V_E = V_G
h^{-2}$.

The additive genetic variance contributed by a single locus within
population~$i$ is $2p_i(1-p_i)\alpha^2$, where $\alpha=2$ is the
average effect of a gene substitution \citep{Fisher:AE-11-53}. For
either population, this equals $2(1 - 4\delta^2)$. Summing across
loci,
\[
V_G = 2K(1 - 4\delta^2)
\]
The mean difference between the character values of the two
populations, expressed in units of the within-population standard
deviation, is
\begin{eqnarray*}
  \frac{D}{\sqrt{V_G + V_E}} &=&
  \frac{8K\delta}{\sqrt{h^{-2}2K(1 - 4\delta^2)}}\\
  &\approx& 5.7 h\sqrt{K}\delta
\end{eqnarray*}
To make this concrete, let us assume that $\delta=0.1$ so that
$F_{ST}\approx 0.02$---a value much smaller than that of the
continental human populations. If heritability is $h^2=0.5$, and the
number of loci is $K=100$, then the two populations differ by 4
standard deviations. If instead there are 1000 loci, the difference
exceeds 12 standard deviations. This shows that small differences in
allele frequencies can cause large population differences in the
quantitative character they control.

\citet{Lewontin:EB-6-381}

\printbibliography

\end{document}

